% paper83-copilot-optimization-advisor.tex — Copilot Optimization Advisor
%   Advisor for Suggesting and Validating Optimization Strategies
%   Paper 83 of the Judgment Geometry series.
% Compile: pdflatex paper83-copilot-optimization-advisor
\documentclass[11pt]{article}
\usepackage{lmodern}
% jugeo-common.tex — shared preamble for all Judgment Geometry papers
% Include via: % jugeo-common.tex — shared preamble for all Judgment Geometry papers
% Include via: % jugeo-common.tex — shared preamble for all Judgment Geometry papers
% Include via: \input{jugeo-common}

% ── Packages ────────────────────────────────────────────────────────────
\usepackage{amsmath,amssymb,amsthm,mathtools}
\usepackage{tikz-cd}
\usepackage{listings}
\usepackage{xcolor}
\usepackage{xspace}
\usepackage{booktabs}
\usepackage{multirow}
\usepackage{enumitem}
\usepackage[expansion=false]{microtype}
\usepackage{hyperref}
\usepackage{cleveref}
% \usepackage{thmtools}  % disabled: not available in this TeX installation
\usepackage{float}
\usepackage{caption}
\usepackage{subcaption}
\usepackage{tabularx}
\usepackage{stmaryrd}       % \llbracket, \rrbracket
\usepackage{bbm}            % \mathbbm{1}

% ── Page geometry ───────────────────────────────────────────────────────
\usepackage[margin=1in]{geometry}

% ── Theorem environments ────────────────────────────────────────────────
\theoremstyle{plain}
\newtheorem{theorem}{Theorem}[section]
\newtheorem{lemma}[theorem]{Lemma}
\newtheorem{proposition}[theorem]{Proposition}
\newtheorem{corollary}[theorem]{Corollary}

\theoremstyle{definition}
\newtheorem{definition}[theorem]{Definition}
\newtheorem{example}[theorem]{Example}
\newtheorem{construction}[theorem]{Construction}

\theoremstyle{remark}
\newtheorem{remark}[theorem]{Remark}
\newtheorem{notation}[theorem]{Notation}

% ── Python listings style ──────────────────────────────────────────────
\definecolor{jugeoBlue}{HTML}{2563EB}
\definecolor{jugeoGreen}{HTML}{059669}
\definecolor{jugeoRed}{HTML}{DC2626}
\definecolor{jugeoPurple}{HTML}{7C3AED}
\definecolor{jugeoGray}{HTML}{6B7280}
\definecolor{jugeoBg}{HTML}{F8FAFC}

\lstdefinestyle{jugeo-python}{
  language=Python,
  basicstyle=\ttfamily\small,
  keywordstyle=\color{jugeoBlue}\bfseries,
  stringstyle=\color{jugeoGreen},
  commentstyle=\color{jugeoGray}\itshape,
  numberstyle=\tiny\color{jugeoGray},
  backgroundcolor=\color{jugeoBg},
  numbers=left,
  numbersep=6pt,
  frame=single,
  framerule=0.4pt,
  rulecolor=\color{jugeoGray!40},
  breaklines=true,
  breakatwhitespace=true,
  showstringspaces=false,
  tabsize=4,
  xleftmargin=1.5em,
  framexleftmargin=1.5em,
  morekeywords={dataclass,frozen,field,Enum,IntEnum,auto,Any,
                Mapping,Sequence,Optional,match,case},
  literate={→}{{$\to$}}1 {←}{{$\leftarrow$}}1
           {≤}{{$\leq$}}1 {≥}{{$\geq$}}1 {≠}{{$\neq$}}1
           {∈}{{$\in$}}1 {∀}{{$\forall$}}1 {∃}{{$\exists$}}1
           {⊕}{{$\oplus$}}1 {⊖}{{$\ominus$}}1,
  escapeinside={(*@}{@*)},
}
\lstset{style=jugeo-python}

\lstdefinestyle{jugeo-output}{
  basicstyle=\ttfamily\small,
  backgroundcolor=\color{jugeoBg},
  frame=single,
  framerule=0.4pt,
  rulecolor=\color{jugeoGray!40},
  breaklines=true,
  showstringspaces=false,
  xleftmargin=1.5em,
  framexleftmargin=0.5em,
  numbers=none,
}

% ── JuGeo-specific macros ──────────────────────────────────────────────

% Judgment 8-tuple
\newcommand{\J}{\mathcal{J}}
\newcommand{\Jud}[8]{(#1,\,#2,\,#3,\,#4,\,#5,\,#6,\,#7,\,#8)}
\newcommand{\JudShort}[1]{\J_{#1}}

% Components
\newcommand{\coord}{\mathsf{c}}            % coordinate
\newcommand{\prop}{\varphi}                 % proposition
\newcommand{\carrier}{\mathcal{A}}          % carrier type
\newcommand{\evid}{\mathcal{E}}             % evidence bundle
\newcommand{\oblig}{\mathcal{O}}            % obligations
\newcommand{\obstr}{\mathcal{B}}            % obstructions
\newcommand{\trust}{\mathcal{T}}            % trust annotation
\newcommand{\prov}{\Pi}                     % provenance

% Trust algebra
\newcommand{\Talg}{\mathfrak{T}}            % trust ordered algebra
\newcommand{\Eadm}{\mathcal{E}_{\mathrm{adm}}}
\newcommand{\tleq}{\preceq}                 % trust ordering
\newcommand{\tjoin}{\oplus}                 % conservative join
\newcommand{\tatten}{\ominus}               % attenuation
\newcommand{\tpromote}{\uparrow_\pi}        % promotion
\newcommand{\tdemote}{\downarrow_\chi}       % demotion

% Trust tiers
\newcommand{\Tcontra}{\ensuremath{\mathsf{CONTRADICTED}}}
\newcommand{\Tunver}{\ensuremath{\mathsf{UNVERIFIED}}}
\newcommand{\Tcopilot}{\ensuremath{\mathsf{COPILOT}}}
\newcommand{\Toracle}{\ensuremath{\mathsf{ORACLE}}}
\newcommand{\Truntime}{\ensuremath{\mathsf{RUNTIME}}}
\newcommand{\Tsolver}{\ensuremath{\mathsf{SOLVER}}}
\newcommand{\Tproof}{\ensuremath{\mathsf{PROOF}}}

% Site and geometry
\newcommand{\Site}{\mathcal{S}}             % semantic site
\newcommand{\Cov}{\mathrm{Cov}}            % covering family
\newcommand{\Ob}{\mathrm{Ob}}              % objects
\newcommand{\Mor}{\mathrm{Mor}}            % morphisms
\newcommand{\res}{\mathrm{res}}            % restriction
\newcommand{\Cech}{\check{C}}              % Čech complex
\newcommand{\Hcoh}[1]{H^{#1}}             % cohomology group

% Descent
\newcommand{\descent}{\mathrm{desc}}
\newcommand{\glue}{\mathrm{glue}}
\newcommand{\overlap}[2]{#1 \cap #2}

% Semantic moves
\newcommand{\VERIFY}{\textsc{Verify}}
\newcommand{\CONSTRUCT}{\textsc{Construct}}
\newcommand{\NORMALIZE}{\textsc{Normalize}}
\newcommand{\GLUE}{\textsc{Glue}}
\newcommand{\RESTRICT}{\textsc{Restrict}}
\newcommand{\TRANSPORT}{\textsc{Transport}}
\newcommand{\REPAIR}{\textsc{Repair}}
\newcommand{\PROMOTE}{\textsc{Promote}}
\newcommand{\CHALLENGE}{\textsc{Challenge}}
\newcommand{\ESCALATE}{\textsc{Escalate}}

% Fragments
\newcommand{\QFLIA}{\texttt{QF\_LIA}}
\newcommand{\QFLRA}{\texttt{QF\_LRA}}
\newcommand{\QFBV}{\texttt{QF\_BV}}
\newcommand{\QFUF}{\texttt{QF\_UF}}

% Misc
\newcommand{\Python}{\textsc{Python}}
\newcommand{\JuGeo}{\textsc{JuGeo}}
\newcommand{\LEAN}{\textsc{Lean}\,4}
\newcommand{\Fstar}{F$^{\star}$}
\newcommand{\Coq}{\textsc{Coq}}
\newcommand{\Dafny}{\textsc{Dafny}}

% Common shortcuts
\newcommand{\ie}{i.e.\@\xspace}
\newcommand{\eg}{e.g.\@\xspace}
\newcommand{\cf}{cf.\@\xspace}
\newcommand{\wrt}{w.r.t.\@\xspace}

% Evaluation shortcuts
\newcommand{\numcases}{300}
\newcommand{\numspec}{100}
\newcommand{\numeq}{100}
\newcommand{\numbug}{100}
\newcommand{\medtime}{6.5\,ms}
\newcommand{\totaltime}{1.949\,s}


% ── Packages ────────────────────────────────────────────────────────────
\usepackage{amsmath,amssymb,amsthm,mathtools}
\usepackage{tikz-cd}
\usepackage{listings}
\usepackage{xcolor}
\usepackage{xspace}
\usepackage{booktabs}
\usepackage{multirow}
\usepackage{enumitem}
\usepackage[expansion=false]{microtype}
\usepackage{hyperref}
\usepackage{cleveref}
% \usepackage{thmtools}  % disabled: not available in this TeX installation
\usepackage{float}
\usepackage{caption}
\usepackage{subcaption}
\usepackage{tabularx}
\usepackage{stmaryrd}       % \llbracket, \rrbracket
\usepackage{bbm}            % \mathbbm{1}

% ── Page geometry ───────────────────────────────────────────────────────
\usepackage[margin=1in]{geometry}

% ── Theorem environments ────────────────────────────────────────────────
\theoremstyle{plain}
\newtheorem{theorem}{Theorem}[section]
\newtheorem{lemma}[theorem]{Lemma}
\newtheorem{proposition}[theorem]{Proposition}
\newtheorem{corollary}[theorem]{Corollary}

\theoremstyle{definition}
\newtheorem{definition}[theorem]{Definition}
\newtheorem{example}[theorem]{Example}
\newtheorem{construction}[theorem]{Construction}

\theoremstyle{remark}
\newtheorem{remark}[theorem]{Remark}
\newtheorem{notation}[theorem]{Notation}

% ── Python listings style ──────────────────────────────────────────────
\definecolor{jugeoBlue}{HTML}{2563EB}
\definecolor{jugeoGreen}{HTML}{059669}
\definecolor{jugeoRed}{HTML}{DC2626}
\definecolor{jugeoPurple}{HTML}{7C3AED}
\definecolor{jugeoGray}{HTML}{6B7280}
\definecolor{jugeoBg}{HTML}{F8FAFC}

\lstdefinestyle{jugeo-python}{
  language=Python,
  basicstyle=\ttfamily\small,
  keywordstyle=\color{jugeoBlue}\bfseries,
  stringstyle=\color{jugeoGreen},
  commentstyle=\color{jugeoGray}\itshape,
  numberstyle=\tiny\color{jugeoGray},
  backgroundcolor=\color{jugeoBg},
  numbers=left,
  numbersep=6pt,
  frame=single,
  framerule=0.4pt,
  rulecolor=\color{jugeoGray!40},
  breaklines=true,
  breakatwhitespace=true,
  showstringspaces=false,
  tabsize=4,
  xleftmargin=1.5em,
  framexleftmargin=1.5em,
  morekeywords={dataclass,frozen,field,Enum,IntEnum,auto,Any,
                Mapping,Sequence,Optional,match,case},
  literate={→}{{$\to$}}1 {←}{{$\leftarrow$}}1
           {≤}{{$\leq$}}1 {≥}{{$\geq$}}1 {≠}{{$\neq$}}1
           {∈}{{$\in$}}1 {∀}{{$\forall$}}1 {∃}{{$\exists$}}1
           {⊕}{{$\oplus$}}1 {⊖}{{$\ominus$}}1,
  escapeinside={(*@}{@*)},
}
\lstset{style=jugeo-python}

\lstdefinestyle{jugeo-output}{
  basicstyle=\ttfamily\small,
  backgroundcolor=\color{jugeoBg},
  frame=single,
  framerule=0.4pt,
  rulecolor=\color{jugeoGray!40},
  breaklines=true,
  showstringspaces=false,
  xleftmargin=1.5em,
  framexleftmargin=0.5em,
  numbers=none,
}

% ── JuGeo-specific macros ──────────────────────────────────────────────

% Judgment 8-tuple
\newcommand{\J}{\mathcal{J}}
\newcommand{\Jud}[8]{(#1,\,#2,\,#3,\,#4,\,#5,\,#6,\,#7,\,#8)}
\newcommand{\JudShort}[1]{\J_{#1}}

% Components
\newcommand{\coord}{\mathsf{c}}            % coordinate
\newcommand{\prop}{\varphi}                 % proposition
\newcommand{\carrier}{\mathcal{A}}          % carrier type
\newcommand{\evid}{\mathcal{E}}             % evidence bundle
\newcommand{\oblig}{\mathcal{O}}            % obligations
\newcommand{\obstr}{\mathcal{B}}            % obstructions
\newcommand{\trust}{\mathcal{T}}            % trust annotation
\newcommand{\prov}{\Pi}                     % provenance

% Trust algebra
\newcommand{\Talg}{\mathfrak{T}}            % trust ordered algebra
\newcommand{\Eadm}{\mathcal{E}_{\mathrm{adm}}}
\newcommand{\tleq}{\preceq}                 % trust ordering
\newcommand{\tjoin}{\oplus}                 % conservative join
\newcommand{\tatten}{\ominus}               % attenuation
\newcommand{\tpromote}{\uparrow_\pi}        % promotion
\newcommand{\tdemote}{\downarrow_\chi}       % demotion

% Trust tiers
\newcommand{\Tcontra}{\ensuremath{\mathsf{CONTRADICTED}}}
\newcommand{\Tunver}{\ensuremath{\mathsf{UNVERIFIED}}}
\newcommand{\Tcopilot}{\ensuremath{\mathsf{COPILOT}}}
\newcommand{\Toracle}{\ensuremath{\mathsf{ORACLE}}}
\newcommand{\Truntime}{\ensuremath{\mathsf{RUNTIME}}}
\newcommand{\Tsolver}{\ensuremath{\mathsf{SOLVER}}}
\newcommand{\Tproof}{\ensuremath{\mathsf{PROOF}}}

% Site and geometry
\newcommand{\Site}{\mathcal{S}}             % semantic site
\newcommand{\Cov}{\mathrm{Cov}}            % covering family
\newcommand{\Ob}{\mathrm{Ob}}              % objects
\newcommand{\Mor}{\mathrm{Mor}}            % morphisms
\newcommand{\res}{\mathrm{res}}            % restriction
\newcommand{\Cech}{\check{C}}              % Čech complex
\newcommand{\Hcoh}[1]{H^{#1}}             % cohomology group

% Descent
\newcommand{\descent}{\mathrm{desc}}
\newcommand{\glue}{\mathrm{glue}}
\newcommand{\overlap}[2]{#1 \cap #2}

% Semantic moves
\newcommand{\VERIFY}{\textsc{Verify}}
\newcommand{\CONSTRUCT}{\textsc{Construct}}
\newcommand{\NORMALIZE}{\textsc{Normalize}}
\newcommand{\GLUE}{\textsc{Glue}}
\newcommand{\RESTRICT}{\textsc{Restrict}}
\newcommand{\TRANSPORT}{\textsc{Transport}}
\newcommand{\REPAIR}{\textsc{Repair}}
\newcommand{\PROMOTE}{\textsc{Promote}}
\newcommand{\CHALLENGE}{\textsc{Challenge}}
\newcommand{\ESCALATE}{\textsc{Escalate}}

% Fragments
\newcommand{\QFLIA}{\texttt{QF\_LIA}}
\newcommand{\QFLRA}{\texttt{QF\_LRA}}
\newcommand{\QFBV}{\texttt{QF\_BV}}
\newcommand{\QFUF}{\texttt{QF\_UF}}

% Misc
\newcommand{\Python}{\textsc{Python}}
\newcommand{\JuGeo}{\textsc{JuGeo}}
\newcommand{\LEAN}{\textsc{Lean}\,4}
\newcommand{\Fstar}{F$^{\star}$}
\newcommand{\Coq}{\textsc{Coq}}
\newcommand{\Dafny}{\textsc{Dafny}}

% Common shortcuts
\newcommand{\ie}{i.e.\@\xspace}
\newcommand{\eg}{e.g.\@\xspace}
\newcommand{\cf}{cf.\@\xspace}
\newcommand{\wrt}{w.r.t.\@\xspace}

% Evaluation shortcuts
\newcommand{\numcases}{300}
\newcommand{\numspec}{100}
\newcommand{\numeq}{100}
\newcommand{\numbug}{100}
\newcommand{\medtime}{6.5\,ms}
\newcommand{\totaltime}{1.949\,s}


% ── Packages ────────────────────────────────────────────────────────────
\usepackage{amsmath,amssymb,amsthm,mathtools}
\usepackage{tikz-cd}
\usepackage{listings}
\usepackage{xcolor}
\usepackage{xspace}
\usepackage{booktabs}
\usepackage{multirow}
\usepackage{enumitem}
\usepackage[expansion=false]{microtype}
\usepackage{hyperref}
\usepackage{cleveref}
% \usepackage{thmtools}  % disabled: not available in this TeX installation
\usepackage{float}
\usepackage{caption}
\usepackage{subcaption}
\usepackage{tabularx}
\usepackage{stmaryrd}       % \llbracket, \rrbracket
\usepackage{bbm}            % \mathbbm{1}

% ── Page geometry ───────────────────────────────────────────────────────
\usepackage[margin=1in]{geometry}

% ── Theorem environments ────────────────────────────────────────────────
\theoremstyle{plain}
\newtheorem{theorem}{Theorem}[section]
\newtheorem{lemma}[theorem]{Lemma}
\newtheorem{proposition}[theorem]{Proposition}
\newtheorem{corollary}[theorem]{Corollary}

\theoremstyle{definition}
\newtheorem{definition}[theorem]{Definition}
\newtheorem{example}[theorem]{Example}
\newtheorem{construction}[theorem]{Construction}

\theoremstyle{remark}
\newtheorem{remark}[theorem]{Remark}
\newtheorem{notation}[theorem]{Notation}

% ── Python listings style ──────────────────────────────────────────────
\definecolor{jugeoBlue}{HTML}{2563EB}
\definecolor{jugeoGreen}{HTML}{059669}
\definecolor{jugeoRed}{HTML}{DC2626}
\definecolor{jugeoPurple}{HTML}{7C3AED}
\definecolor{jugeoGray}{HTML}{6B7280}
\definecolor{jugeoBg}{HTML}{F8FAFC}

\lstdefinestyle{jugeo-python}{
  language=Python,
  basicstyle=\ttfamily\small,
  keywordstyle=\color{jugeoBlue}\bfseries,
  stringstyle=\color{jugeoGreen},
  commentstyle=\color{jugeoGray}\itshape,
  numberstyle=\tiny\color{jugeoGray},
  backgroundcolor=\color{jugeoBg},
  numbers=left,
  numbersep=6pt,
  frame=single,
  framerule=0.4pt,
  rulecolor=\color{jugeoGray!40},
  breaklines=true,
  breakatwhitespace=true,
  showstringspaces=false,
  tabsize=4,
  xleftmargin=1.5em,
  framexleftmargin=1.5em,
  morekeywords={dataclass,frozen,field,Enum,IntEnum,auto,Any,
                Mapping,Sequence,Optional,match,case},
  literate={→}{{$\to$}}1 {←}{{$\leftarrow$}}1
           {≤}{{$\leq$}}1 {≥}{{$\geq$}}1 {≠}{{$\neq$}}1
           {∈}{{$\in$}}1 {∀}{{$\forall$}}1 {∃}{{$\exists$}}1
           {⊕}{{$\oplus$}}1 {⊖}{{$\ominus$}}1,
  escapeinside={(*@}{@*)},
}
\lstset{style=jugeo-python}

\lstdefinestyle{jugeo-output}{
  basicstyle=\ttfamily\small,
  backgroundcolor=\color{jugeoBg},
  frame=single,
  framerule=0.4pt,
  rulecolor=\color{jugeoGray!40},
  breaklines=true,
  showstringspaces=false,
  xleftmargin=1.5em,
  framexleftmargin=0.5em,
  numbers=none,
}

% ── JuGeo-specific macros ──────────────────────────────────────────────

% Judgment 8-tuple
\newcommand{\J}{\mathcal{J}}
\newcommand{\Jud}[8]{(#1,\,#2,\,#3,\,#4,\,#5,\,#6,\,#7,\,#8)}
\newcommand{\JudShort}[1]{\J_{#1}}

% Components
\newcommand{\coord}{\mathsf{c}}            % coordinate
\newcommand{\prop}{\varphi}                 % proposition
\newcommand{\carrier}{\mathcal{A}}          % carrier type
\newcommand{\evid}{\mathcal{E}}             % evidence bundle
\newcommand{\oblig}{\mathcal{O}}            % obligations
\newcommand{\obstr}{\mathcal{B}}            % obstructions
\newcommand{\trust}{\mathcal{T}}            % trust annotation
\newcommand{\prov}{\Pi}                     % provenance

% Trust algebra
\newcommand{\Talg}{\mathfrak{T}}            % trust ordered algebra
\newcommand{\Eadm}{\mathcal{E}_{\mathrm{adm}}}
\newcommand{\tleq}{\preceq}                 % trust ordering
\newcommand{\tjoin}{\oplus}                 % conservative join
\newcommand{\tatten}{\ominus}               % attenuation
\newcommand{\tpromote}{\uparrow_\pi}        % promotion
\newcommand{\tdemote}{\downarrow_\chi}       % demotion

% Trust tiers
\newcommand{\Tcontra}{\ensuremath{\mathsf{CONTRADICTED}}}
\newcommand{\Tunver}{\ensuremath{\mathsf{UNVERIFIED}}}
\newcommand{\Tcopilot}{\ensuremath{\mathsf{COPILOT}}}
\newcommand{\Toracle}{\ensuremath{\mathsf{ORACLE}}}
\newcommand{\Truntime}{\ensuremath{\mathsf{RUNTIME}}}
\newcommand{\Tsolver}{\ensuremath{\mathsf{SOLVER}}}
\newcommand{\Tproof}{\ensuremath{\mathsf{PROOF}}}

% Site and geometry
\newcommand{\Site}{\mathcal{S}}             % semantic site
\newcommand{\Cov}{\mathrm{Cov}}            % covering family
\newcommand{\Ob}{\mathrm{Ob}}              % objects
\newcommand{\Mor}{\mathrm{Mor}}            % morphisms
\newcommand{\res}{\mathrm{res}}            % restriction
\newcommand{\Cech}{\check{C}}              % Čech complex
\newcommand{\Hcoh}[1]{H^{#1}}             % cohomology group

% Descent
\newcommand{\descent}{\mathrm{desc}}
\newcommand{\glue}{\mathrm{glue}}
\newcommand{\overlap}[2]{#1 \cap #2}

% Semantic moves
\newcommand{\VERIFY}{\textsc{Verify}}
\newcommand{\CONSTRUCT}{\textsc{Construct}}
\newcommand{\NORMALIZE}{\textsc{Normalize}}
\newcommand{\GLUE}{\textsc{Glue}}
\newcommand{\RESTRICT}{\textsc{Restrict}}
\newcommand{\TRANSPORT}{\textsc{Transport}}
\newcommand{\REPAIR}{\textsc{Repair}}
\newcommand{\PROMOTE}{\textsc{Promote}}
\newcommand{\CHALLENGE}{\textsc{Challenge}}
\newcommand{\ESCALATE}{\textsc{Escalate}}

% Fragments
\newcommand{\QFLIA}{\texttt{QF\_LIA}}
\newcommand{\QFLRA}{\texttt{QF\_LRA}}
\newcommand{\QFBV}{\texttt{QF\_BV}}
\newcommand{\QFUF}{\texttt{QF\_UF}}

% Misc
\newcommand{\Python}{\textsc{Python}}
\newcommand{\JuGeo}{\textsc{JuGeo}}
\newcommand{\LEAN}{\textsc{Lean}\,4}
\newcommand{\Fstar}{F$^{\star}$}
\newcommand{\Coq}{\textsc{Coq}}
\newcommand{\Dafny}{\textsc{Dafny}}

% Common shortcuts
\newcommand{\ie}{i.e.\@\xspace}
\newcommand{\eg}{e.g.\@\xspace}
\newcommand{\cf}{cf.\@\xspace}
\newcommand{\wrt}{w.r.t.\@\xspace}

% Evaluation shortcuts
\newcommand{\numcases}{300}
\newcommand{\numspec}{100}
\newcommand{\numeq}{100}
\newcommand{\numbug}{100}
\newcommand{\medtime}{6.5\,ms}
\newcommand{\totaltime}{1.949\,s}


% ——— Paper-specific macros ———
\newcommand{\OptAdv}{\mathsf{CopilotOptimizationAdvisor}}
\newcommand{\OptStrat}{\mathrm{optimizationStrategy}}
\newcommand{\PerfPred}{\mathrm{performancePredict}}
\newcommand{\NonReg}{\mathrm{nonRegression}}
\newcommand{\OptValid}{\mathrm{optimizationValid}}
\newcommand{\TrustPres}{\mathrm{trustPreserve}}
\newcommand{\SearchRed}{\mathrm{searchSpaceReduce}}
\newcommand{\ValComp}{\mathrm{validationComplete}}
\newcommand{\CopChan}{\mathsf{CopilotChannel}}
\newcommand{\ChanMon}{\mathsf{ChannelMonitor}}
\newcommand{\TrustCeil}{\mathrm{ceiling}}
\newcommand{\HealthInd}{\mathsf{HealthIndicator}}
\newcommand{\JuGeoAPI}{\mathsf{JuGeoAPI}}

\title{\textbf{Copilot Optimization Advisor:\\ Advisor for Suggesting and Validating Optimization Strategies}}
\author{Halley Young}
\date{Paper~83 --- Judgment Geometry Series}

\begin{document}
\maketitle

% ——— Abstract ———
\begin{abstract}
We present the $\OptAdv$, a Copilot-guided advisor for suggesting and validating optimization strategies in the \JuGeo{} verification framework. The advisor proposes optimizations, predicts performance, and checks for non-regression, all within a trust-calibrated evidence architecture. We formalize the validation workflow, prove soundness, non-regression, search space reduction, trust preservation, and validation completeness, and provide full proofs. Code listings demonstrate $\CopChan$ for strategy suggestion, $\ChanMon$ for performance tracking, $\JuGeoAPI$ for verification, and $\HealthInd$ for subsystem health. The appendix contains additional proofs and technical details.
\end{abstract}

% ============================================================
%  Section 1 — Introduction
% ============================================================
\section{Introduction}\label{sec:intro}

Verification tasks in the \JuGeo{} framework involve searching
large proof spaces under tight resource budgets.  Choosing
the right verification strategy---which lemmas to inline,
which sub-goals to split, which channel to route evidence
requests through---can reduce verification time by orders
of magnitude.  Yet the space of possible optimisations is
itself large, and a na\"{\i}ve search over strategies is
often as expensive as the verification it aims to accelerate.

This paper introduces the $\OptAdv$, a Copilot-guided
advisor that suggests optimisation strategies and validates
them against a correctness preservation invariant.  The
advisor operates at the $\Tcopilot$ trust tier: its
suggestions are proposals that the kernel may accept or
reject.  Validation is delegated to higher-trust channels
(solver, runtime) that independently confirm that the
optimised verification yields the same result as the
original.

\paragraph{Contributions.}
\begin{enumerate}
\item A formal model of verification optimisation as a
      rewriting system on proof obligations (\S\ref{sec:model}).
\item The design of $\OptAdv$ with three phases:
      strategy suggestion, performance prediction, and
      non-regression checking (\S\ref{sec:advisor}).
\item A validation framework that uses the solver channel
      to confirm correctness preservation (\S\ref{sec:validation}).
\item Five theorems: optimisation soundness, non-regression,
      search-space reduction, trust preservation, and
      validation completeness (\S\ref{sec:formal}).
\item Full \Python{} code illustrating integration with
      $\CopChan$, $\ChanMon$, $\JuGeoAPI$, and
      $\HealthInd$ (\S\ref{sec:code}).
\end{enumerate}

\paragraph{Paper outline.}
Section~\ref{sec:background} recalls background material.
Section~\ref{sec:model} formalises verification optimisation.
Section~\ref{sec:advisor} presents the advisor design.
Section~\ref{sec:validation} describes validation.
Section~\ref{sec:code} gives implementation details.
Section~\ref{sec:formal} contains the main theorems.
Section~\ref{sec:discussion} discusses limitations.
Section~\ref{sec:related} surveys related work.
Section~\ref{sec:conclusion} concludes.
Appendix~\ref{app:proofs} collects auxiliary proofs.

% ============================================================
%  Section 2 — Background
% ============================================================
\section{Background}\label{sec:background}

\subsection{Trust lattice recap}

Recall the trust ordering from earlier papers:
\[
  \Tcontra \;\tleq\; \Tunver \;\tleq\; \Tcopilot
  \;\tleq\; \Toracle \;\tleq\; \Truntime
  \;\tleq\; \Tsolver \;\tleq\; \Tproof.
\]
The optimisation advisor operates at $\Tcopilot$; its
suggestions are not trusted until validated by a
higher-trust channel.

\subsection{Evidence channels and monitoring}

Evidence channels are the conduits through which the
kernel obtains evidence.  The \texttt{ChannelMonitor}
records request--response pairs and computes summary
statistics: success rate, latency percentiles, and
throughput.  These statistics are the raw material
for the optimisation advisor's performance predictions.

\subsection{Proof obligations and rewriting}

A \emph{proof obligation} is a pair $(c, \varphi)$
where $c$ is a coordinate in the proof tree and
$\varphi$ is a proposition to be verified.  A
\emph{verification strategy} is a sequence of
rewriting steps that transforms a proof obligation
into a (possibly empty) set of simpler obligations.

\begin{definition}[Rewriting system]\label{def:rewrite}
A \emph{verification rewriting system} is a pair
$(\mathcal{O}, \Rightarrow)$ where $\mathcal{O}$
is the set of proof obligations and
$\Rightarrow \;\subseteq \mathcal{O} \times
\mathcal{P}(\mathcal{O})$ is the rewriting relation.
We write $o \Rightarrow O'$ to mean that obligation~$o$
can be rewritten into the set $O'$ of sub-obligations.
\end{definition}

\begin{definition}[Optimisation]\label{def:optimisation}
An \emph{optimisation} is a pair $(\rho, \pi)$ where
$\rho \colon \mathcal{O} \to \mathcal{O}$ is a
transformation on obligations and $\pi$ is a proof
that $\rho$ preserves validity:
$\mathrm{valid}(o) \iff \mathrm{valid}(\rho(o))$.
\end{definition}

% ============================================================
%  Section 3 — The Optimisation Advisor
% ============================================================
\section{The Optimisation Advisor}\label{sec:advisor}

\subsection{Strategy suggestion}

The $\OptAdv$ receives the current proof obligation,
the rewriting history, and the monitor statistics.
It queries the Copilot channel for a ranked list of
optimisation strategies.

\begin{definition}[Strategy space]\label{def:stratspace}
The \emph{strategy space} $\mathcal{S}$ is the set of
all valid optimisations.  A strategy $s \in \mathcal{S}$
is a named rewriting rule together with estimated cost
and benefit metrics:
\[
  s = (\mathsf{name}, \rho, \mathsf{cost}, \mathsf{benefit}).
\]
The \emph{net benefit} is $\mathsf{benefit} - \mathsf{cost}$.
\end{definition}

\begin{definition}[Suggestion ranking]\label{def:sugrank}
The advisor ranks strategies by expected net benefit.
Given monitor data~$d$ and obligation~$o$, the ranking
function is:
\[
  \mathrm{rank}(s, o, d) =
    \frac{\mathsf{benefit}(s, o)}{\mathsf{cost}(s, o)}
    \times \mathrm{success\_rate}(d, s.\mathsf{name}).
\]
\end{definition}

\subsection{Performance prediction}

Before suggesting an optimisation, the advisor predicts
its effect on verification time.  The prediction model
uses the monitor's latency percentiles and success rates.

\begin{definition}[Performance model]\label{def:perfmodel}
A \emph{performance model} is a function
$M \colon \mathcal{S} \times \mathcal{O} \times
\mathsf{MonitorData} \to \mathbb{R}_{\geq 0}$
that estimates the verification time under strategy~$s$
for obligation~$o$ given historical data~$d$.  We require:
\begin{enumerate}
\item $M$ is monotone in cost: if $s$ has higher cost
      than $s'$, then $M(s, o, d) \geq M(s', o, d)$.
\item $M$ is calibrated: the mean prediction error
      over the last $n$ verifications is below a
      threshold~$\epsilon$.
\end{enumerate}
\end{definition}

\subsection{Non-regression checking}

After suggesting an optimisation, the advisor checks
that the optimised obligation has the same validity
status as the original.  This check is delegated to
the solver channel.

\begin{definition}[Non-regression check]\label{def:nonreg}
A \emph{non-regression check} for optimisation~$\rho$
on obligation~$o$ is a verification that:
\[
  \mathrm{verify}(\rho(o)) = \mathrm{verify}(o).
\]
The check succeeds if both sides agree (both valid or
both invalid).  It fails if the optimisation changes
the verification outcome.
\end{definition}

\begin{remark}
Non-regression checking is the key safety mechanism.
Even though the Copilot channel suggests optimisations
at $\Tcopilot$, the solver channel validates them at
$\Tsolver$.  This two-phase approach prevents unsound
optimisations from being applied.
\end{remark}

% ============================================================
%  Section 4 — Validation Framework
% ============================================================
\section{Optimisation Validation}\label{sec:validation}

\subsection{Correctness preservation}

\begin{definition}[Correctness-preserving optimisation]%
\label{def:corrpres}
An optimisation $\rho$ is \emph{correctness-preserving}
if for every obligation~$o$:
\[
  \mathrm{valid}(o) \iff \mathrm{valid}(\rho(o)).
\]
Equivalently, $\rho$ is a validity-preserving endomorphism
on $(\mathcal{O}, \mathrm{valid})$.
\end{definition}

\subsection{Solver-checked validation}

The validation protocol proceeds as follows:
\begin{enumerate}
\item The advisor suggests optimisation~$\rho$ for obligation~$o$.
\item The kernel submits both $o$ and $\rho(o)$ to the solver
      channel.
\item The solver returns verdicts $v$ and $v'$.
\item If $v = v'$, the optimisation is accepted.
\item If $v \neq v'$, the optimisation is rejected.
\end{enumerate}

\begin{definition}[Validation oracle]\label{def:valoracle}
A \emph{validation oracle} is a function
$V \colon \mathcal{O} \times \mathcal{O} \to \{
\mathsf{accept}, \mathsf{reject}, \mathsf{unknown}\}$
defined by:
\[
  V(o, o') = \begin{cases}
    \mathsf{accept} & \text{if } \mathrm{verify}(o) =
      \mathrm{verify}(o'), \\
    \mathsf{reject} & \text{if } \mathrm{verify}(o) \neq
      \mathrm{verify}(o'), \\
    \mathsf{unknown} & \text{if either verification times out.}
  \end{cases}
\]
\end{definition}

\subsection{Incremental validation}

For large optimisations that transform many sub-obligations,
incremental validation checks each sub-obligation individually.

\begin{definition}[Incremental validation]\label{def:incval}
Given an optimisation $\rho$ and a set of sub-obligations
$\{o_1, \ldots, o_n\}$, the incremental validator computes
$V(o_i, \rho(o_i))$ for each~$i$ and returns:
\[
  V_{\mathrm{inc}}(\rho, \{o_1, \ldots, o_n\}) =
  \begin{cases}
    \mathsf{accept} & \text{if all } V(o_i, \rho(o_i)) =
      \mathsf{accept}, \\
    \mathsf{reject} & \text{if any } V(o_i, \rho(o_i)) =
      \mathsf{reject}, \\
    \mathsf{unknown} & \text{otherwise.}
  \end{cases}
\]
\end{definition}

\begin{lemma}[Incremental validation soundness]%
\label{lem:incvalsound}
If $V_{\mathrm{inc}}(\rho, O) = \mathsf{accept}$, then
$\rho$ is correctness-preserving on~$O$.
\end{lemma}
\begin{proof}
Each sub-obligation $o_i \in O$ has been individually
validated: $\mathrm{verify}(o_i) = \mathrm{verify}(\rho(o_i))$.
Therefore $\rho$ preserves validity on each element of~$O$.
Since $O$ is the complete set of sub-obligations, $\rho$
is correctness-preserving on~$O$.
\end{proof}

% ============================================================
%  Section 5 — Implementation
% ============================================================
\section{Implementation}\label{sec:code}

\subsection{Strategy suggestion via CopilotChannel}

The following listing shows how the $\OptAdv$ queries
the Copilot channel for optimisation strategies.

\begin{lstlisting}[style=jugeo-python,
  caption={Optimisation strategy suggestion via CopilotChannel},
  label={lst:optstrat}]
from jugeo.evidence.channels import (
    CopilotChannel,
    ChannelConfiguration,
    EvidenceChannel,
    EvidenceRequest,
    EvidenceResponse,
    ChannelJurisdiction,
    ChannelMonitor,
)

class CopilotOptimizationAdvisor:
    """Advisor for suggesting and validating optimisation strategies."""

    def __init__(self, channel: CopilotChannel, monitor: ChannelMonitor):
        self.channel = channel
        self.monitor = monitor
        self.trust_ceiling = channel.TRUST_CEILING
        self.jurisdiction = ChannelJurisdiction.for_channel(
            EvidenceChannel.COPILOT
        )

    def suggest_strategies(self, obligation, rewrite_history, k=5):
        stats = self.monitor.summary()
        latencies = self.monitor.latency_percentiles()
        success = self.monitor.success_rate()

        prompt = self._build_strategy_prompt(
            obligation, rewrite_history, stats, latencies, success
        )
        raw = self.channel.query_llm(prompt, timeout_ms=10000)
        parsed = self.channel.parse_response()
        clamped = self.channel.apply_trust_ceiling()

        strategies = self._rank_strategies(clamped, obligation, stats)
        return strategies[:k]

    def predict_performance(self, strategy, obligation):
        stats = self.monitor.summary()
        latencies = self.monitor.latency_percentiles()
        p50 = latencies.get("p50", 1000)
        p99 = latencies.get("p99", 5000)
        success = self.monitor.success_rate()
        estimated_time = p50 * strategy.get("cost_factor", 1.0)
        return {
            "estimated_ms": estimated_time,
            "confidence": success,
            "p99_bound": p99 * strategy.get("cost_factor", 1.0),
        }

    def _build_strategy_prompt(self, obl, hist, stats, lat, succ):
        return (
            f"Obligation: {obl}\n"
            f"History length: {len(hist)}\n"
            f"Success rate: {succ:.2f}\n"
            f"P50 latency: {lat.get('p50', 'N/A')}ms\n"
            f"Suggest optimization strategies."
        )

    def _rank_strategies(self, strategies, obligation, stats):
        if not isinstance(strategies, list):
            return []
        ranked = sorted(
            strategies,
            key=lambda s: s.get("benefit", 0) / max(s.get("cost", 1), 1),
            reverse=True,
        )
        return ranked
\end{lstlisting}

\subsection{Performance tracking with ChannelMonitor}

The next listing shows how the $\ChanMon$ tracks
performance during optimisation experiments.

\begin{lstlisting}[style=jugeo-python,
  caption={Performance monitoring during optimisation},
  label={lst:perfmon}]
from jugeo.evidence.channels import (
    ChannelMonitor,
    ChannelPool,
    ChannelRouter,
    CopilotChannel,
    SolverChannel,
    ChannelConfiguration,
    EvidenceChannel,
    EvidenceRequest,
    EvidenceResponse,
)

class OptimizationPerformanceTracker:
    """Track performance metrics across optimization experiments."""

    def __init__(self, pool: ChannelPool, monitor: ChannelMonitor):
        self.pool = pool
        self.monitor = monitor
        self.experiments = []

    def run_baseline(self, obligation, channel_name="solver"):
        channel = self.pool.get_channel(channel_name)
        request = EvidenceRequest(
            request_id="opt-baseline-001",
            coordinate=obligation,
            proposition=obligation,
            required_kind=None,
            proposition_kind="verification",
            preferred_channel=EvidenceChannel.SOLVER,
            fallback_channels=[],
            deadline_ms=30000,
            budget=50000,
            metadata={"phase": "baseline"},
        )
        self.monitor.record_request(request)
        response = channel.handle_request(request)
        self.monitor.record_response(response)
        return {
            "request": request,
            "response": response,
            "latency_p50": self.monitor.latency_percentiles().get("p50"),
            "success_rate": self.monitor.success_rate(),
        }

    def run_optimized(self, obligation, strategy, channel_name="solver"):
        optimized_obl = self._apply_strategy(obligation, strategy)
        channel = self.pool.get_channel(channel_name)
        request = EvidenceRequest(
            request_id="opt-test-001",
            coordinate=optimized_obl,
            proposition=optimized_obl,
            required_kind=None,
            proposition_kind="verification",
            preferred_channel=EvidenceChannel.SOLVER,
            fallback_channels=[],
            deadline_ms=30000,
            budget=50000,
            metadata={"phase": "optimized", "strategy": strategy},
        )
        self.monitor.record_request(request)
        response = channel.handle_request(request)
        self.monitor.record_response(response)
        return {
            "request": request,
            "response": response,
            "latency_p50": self.monitor.latency_percentiles().get("p50"),
            "success_rate": self.monitor.success_rate(),
        }

    def compare(self, baseline, optimized):
        baseline_lat = baseline.get("latency_p50", 0)
        opt_lat = optimized.get("latency_p50", 0)
        speedup = baseline_lat / max(opt_lat, 1)
        return {
            "speedup": speedup,
            "baseline_success": baseline.get("success_rate"),
            "optimized_success": optimized.get("success_rate"),
            "summary": self.monitor.summary(),
        }

    def _apply_strategy(self, obligation, strategy):
        return f"optimized({obligation}, {strategy})"
\end{lstlisting}

\subsection{Verification via JuGeoAPI}

The following listing demonstrates using $\JuGeoAPI$
to verify both the original and optimised obligations.

\begin{lstlisting}[style=jugeo-python,
  caption={Non-regression verification via JuGeoAPI},
  label={lst:nonreg}]
from jugeo.interfaces.api import (
    JuGeoAPI,
    OperationKind,
    APIRequest,
    APIResponse,
    APISession,
    APIEventLog,
    APIValidator,
    APIRateLimiter,
    APIAuthenticator,
    CopilotAPIBridge,
    RequestStatus,
)

class NonRegressionChecker:
    """Check that optimisations preserve verification outcomes."""

    def __init__(self, api: JuGeoAPI, event_log: APIEventLog):
        self.api = api
        self.event_log = event_log

    def check(self, original_obligation, optimized_obligation):
        session = self.api.open_session("non-regression")

        orig_request = APIRequest(
            request_id="nonreg-orig-001",
            operation=OperationKind.VERIFY,
            coordinate=original_obligation,
            proposition=original_obligation,
            budget=30000,
            deadline=60000,
            metadata={"phase": "non-regression-original"},
        )
        session.register_request(orig_request)
        orig_result = self.api.verify(
            coordinate=original_obligation,
            proposition=original_obligation,
        )

        opt_request = APIRequest(
            request_id="nonreg-opt-001",
            operation=OperationKind.VERIFY,
            coordinate=optimized_obligation,
            proposition=optimized_obligation,
            budget=30000,
            deadline=60000,
            metadata={"phase": "non-regression-optimized"},
        )
        session.register_request(opt_request)
        opt_result = self.api.verify(
            coordinate=optimized_obligation,
            proposition=optimized_obligation,
        )

        self.event_log.record({
            "original": str(orig_result),
            "optimized": str(opt_result),
        })

        if orig_result.is_successful() == opt_result.is_successful():
            verdict = "accept"
        else:
            verdict = "reject"

        session.checkpoint()
        session.close()
        return {
            "verdict": verdict,
            "original_success": orig_result.is_successful(),
            "optimized_success": opt_result.is_successful(),
        }

    def batch_check(self, pairs):
        results = []
        for orig, opt in pairs:
            result = self.check(orig, opt)
            results.append(result)
        all_accept = all(r["verdict"] == "accept" for r in results)
        return {"all_accept": all_accept, "details": results}
\end{lstlisting}

\subsection{Health monitoring for optimisation subsystem}

The final listing shows how $\HealthInd$ monitors
the health of the optimisation subsystem.

\begin{lstlisting}[style=jugeo-python,
  caption={Health monitoring for the optimisation subsystem},
  label={lst:opthealth}]
from jugeo.kernel.health import (
    HealthStatus,
    HealthDimension,
    HealthIndicator,
    HealthReport,
)
from jugeo.evidence.channels import (
    ChannelPool,
    ChannelMonitor,
    CopilotChannel,
)

class OptimizationHealthMonitor:
    """Monitor health of the optimisation advisory subsystem."""

    def __init__(self, pool: ChannelPool, monitor: ChannelMonitor):
        self.pool = pool
        self.monitor = monitor

    def check_copilot_health(self):
        copilot = self.pool.get_channel("copilot")
        if copilot is None:
            return HealthIndicator(
                dimension=HealthDimension.COPILOT_CONNECTIVITY,
                status=HealthStatus.UNHEALTHY,
                message="Copilot channel not registered",
                value=0.0,
                timestamp=None,
                details={"channel": "copilot"},
            )
        health = copilot.health_check()
        status = (HealthStatus.HEALTHY if health
                  else HealthStatus.DEGRADED)
        return HealthIndicator(
            dimension=HealthDimension.COPILOT_CONNECTIVITY,
            status=status,
            message="Copilot channel operational" if health
                    else "Copilot channel degraded",
            value=1.0 if health else 0.5,
            timestamp=None,
            details={"health_check": health},
        )

    def check_evidence_flow(self):
        success = self.monitor.success_rate()
        if success >= 0.9:
            status = HealthStatus.HEALTHY
        elif success >= 0.5:
            status = HealthStatus.DEGRADED
        else:
            status = HealthStatus.UNHEALTHY
        return HealthIndicator(
            dimension=HealthDimension.EVIDENCE_FLOW,
            status=status,
            message=f"Evidence flow success rate: {success:.2%}",
            value=success,
            timestamp=None,
            details=self.monitor.summary(),
        )

    def full_report(self):
        copilot_ind = self.check_copilot_health()
        flow_ind = self.check_evidence_flow()
        indicators = [copilot_ind, flow_ind]
        if all(i.is_healthy() for i in indicators):
            overall = HealthStatus.HEALTHY
        elif any(i.status == HealthStatus.UNHEALTHY for i in indicators):
            overall = HealthStatus.UNHEALTHY
        else:
            overall = HealthStatus.DEGRADED
        report = HealthReport(
            overall_status=overall,
            indicators=indicators,
            generated_at=None,
            report_id="opt-health-001",
            degradation_reasons=[],
            recovery_suggestions=[],
        )
        return report
\end{lstlisting}

% ============================================================
%  Section 6 — Formal Properties
% ============================================================
\section{Formal Properties}\label{sec:formal}

\begin{theorem}[Optimisation soundness]\label{thm:optsound}
If the advisor suggests optimisation~$\rho$ and the
non-regression check accepts it, then $\rho$ is
correctness-preserving on the tested obligations.
\end{theorem}
\begin{proof}
The non-regression checker (Listing~\ref{lst:nonreg})
submits both the original and optimised obligations
to the solver channel via \texttt{JuGeoAPI.verify()}.
Acceptance requires that both verifications yield the
same outcome (\ie{} both succeed or both fail).
Since the solver channel operates at $\Tsolver$ trust,
its verdicts are reliable up to solver soundness.
Therefore, if the check accepts, then
$\mathrm{verify}(o) = \mathrm{verify}(\rho(o))$,
which by definition means $\rho$ is correctness-preserving
on~$o$.

For the batch case, Lemma~\ref{lem:incvalsound} extends
this to all sub-obligations.
\end{proof}

\begin{theorem}[Non-regression guarantee]\label{thm:nonreg}
No optimisation accepted by the non-regression checker
changes the verification outcome of any tested obligation.
\end{theorem}
\begin{proof}
This is the contrapositive of soundness.  Suppose an
accepted optimisation~$\rho$ changes the outcome of some
obligation~$o$: $\mathrm{verify}(o) \neq
\mathrm{verify}(\rho(o))$.  Then the checker would have
returned ``reject'' for the pair $(o, \rho(o))$.
Contradiction with the assumption that the optimisation
was accepted.
\end{proof}

\begin{theorem}[Search-space reduction]\label{thm:searchred}
If optimisation~$\rho$ reduces the number of sub-obligations
from~$n$ to~$m < n$, then the verification cost decreases
by at least a factor of $n/m$ (assuming linear cost per
sub-obligation).
\end{theorem}
\begin{proof}
Under the linear cost model, verification time is
$T = c \cdot k$ where $c$ is the per-obligation cost and
$k$ is the number of obligations.  The original time is
$T_{\mathrm{orig}} = c \cdot n$ and the optimised time is
$T_{\mathrm{opt}} = c \cdot m$.  The speedup factor is
\[
  \frac{T_{\mathrm{orig}}}{T_{\mathrm{opt}}}
  = \frac{c \cdot n}{c \cdot m}
  = \frac{n}{m}.
\]
Since $m < n$, this factor is greater than~1.

The linear cost assumption holds when all sub-obligations
have comparable complexity.  For heterogeneous obligations,
the actual speedup may differ, but the theorem provides a
lower bound under the stated assumption.
\end{proof}

\begin{theorem}[Trust preservation]\label{thm:trustpres}
The optimisation advisor preserves trust:
$\mathrm{trust}(\OptAdv(\Gamma)) \tleq \Tcopilot$.
\end{theorem}
\begin{proof}
The advisor queries the Copilot channel, which has
trust ceiling $\Tcopilot$.  The \texttt{apply\_trust\_ceiling}
call ensures the response trust does not exceed this
ceiling.  The advisor does not escalate trust:
strategy suggestions inherit the channel's trust.

Note that the non-regression check uses the solver channel
at $\Tsolver$, but this evidence is separate from the
advisor's suggestion.  The suggestion itself remains at
$\Tcopilot$; only the validation result has higher trust.
\end{proof}

\begin{theorem}[Validation completeness]\label{thm:valcomp}
If the solver channel is complete (\ie{} it returns a
definitive verdict for every decidable obligation within
the deadline), then the non-regression checker correctly
classifies every decidable optimisation.
\end{theorem}
\begin{proof}
A decidable obligation is one for which the solver
returns either ``valid'' or ``invalid'' within the
deadline.  Given two decidable obligations $o$ and
$\rho(o)$, the solver returns definitive verdicts
$v$ and $v'$.  The checker compares them:
\begin{itemize}
\item If $v = v'$, the optimisation is accepted (correctly,
      since validity is preserved).
\item If $v \neq v'$, the optimisation is rejected
      (correctly, since validity changed).
\end{itemize}
The only case where the checker returns ``unknown'' is
when one or both verifications time out, which happens
only for undecidable obligations.  Therefore the checker
is complete for decidable obligations.
\end{proof}

% ============================================================
%  Section 7 — Discussion
% ============================================================
\section{Discussion}\label{sec:discussion}

\paragraph{Undecidable obligations.}
When the solver times out, the non-regression checker
returns ``unknown.''  In this case, the advisor falls
back to a conservative strategy: reject the optimisation.
This ensures safety at the cost of missing valid
optimisations for undecidable obligations.

\paragraph{Cost of validation.}
Validating an optimisation requires running the solver
twice (on the original and optimised obligations).
This doubles the solver cost.  In practice, the
optimisation should provide a speedup that more than
compensates for this overhead.  If the predicted
speedup is less than~2, the advisor does not suggest
the optimisation.

\paragraph{Interaction with other advisors.}
The optimisation advisor can be composed with the
research and experiment advisors (Paper~82).  The
research advisor suggests hypotheses, the experiment
advisor designs a verification plan, and the
optimisation advisor suggests how to execute the
plan efficiently.

\paragraph{Stochastic strategies.}
The Copilot channel may suggest different strategies
on different invocations for the same obligation.
The advisor mitigates this by caching responses and
by requiring that any accepted optimisation pass the
deterministic non-regression check.

\paragraph{Health monitoring.}
Listing~\ref{lst:opthealth} shows how the optimisation
subsystem monitors its own health.  If the Copilot
channel is degraded, the advisor reduces its suggestion
rate.  If the evidence flow success rate drops below
50\%, the advisor enters a conservative mode where
only previously validated optimisations are suggested.

% ============================================================
%  Section 8 — Related Work
% ============================================================
\section{Related Work}\label{sec:related}

\paragraph{Automated strategy selection.}
Portfolio solvers~\cite{xu2008satzilla} select among
verification strategies based on problem features.
Our advisor differs in using a language model for
strategy suggestion and a trust lattice for validation.

\paragraph{Compiler optimisation verification.}
Translation validation~\cite{pnueli1998translation}
checks that compiler optimisations preserve semantics.
Our non-regression checker is analogous, applied to
verification strategies rather than compiler passes.

\paragraph{AI for code optimisation.}
Recent work uses language models for code
optimisation~\cite{chen2023codeoptimize}.  We apply
similar ideas to verification optimisation, with the
additional constraint of formal correctness preservation.

\paragraph{Proof search heuristics.}
Ripple~\cite{basin1993rippling}, heuristic
instantiation~\cite{claessen2000quickcheck}, and
learned proof strategies~\cite{bansal2019holist}
all guide proof search.  Our contribution is the
integration with a trust-calibrated advisory framework.

\paragraph{Performance prediction.}
Machine-learning models for runtime
prediction~\cite{hutter2011sequential} and algorithm
selection~\cite{kotthoff2016algorithm} share goals
with our performance model.  We focus on integration
with the \JuGeo{} monitor rather than building a
standalone predictor.

\paragraph{Copilot advisors in \JuGeo{}.}
Papers~80--82 of this series introduce advisors for
heap analysis, scope management, import resolution,
callable analysis, research, and experiment design.
The optimisation advisor complements these by
providing performance-focused guidance.

% ============================================================
%  Section 9 — Conclusion
% ============================================================
\section{Conclusion}\label{sec:conclusion}

We have introduced the $\OptAdv$, a Copilot-guided
advisor that suggests and validates verification
optimisation strategies in the \JuGeo{} framework.
The advisor operates at $\Tcopilot$ trust and relies
on the solver channel for non-regression validation.
Five theorems establish soundness, non-regression,
search-space reduction, trust preservation, and
validation completeness.

Future work includes: (a)~extending the performance
model with learned features from the monitor's time
series, (b)~supporting multi-objective optimisation
(time vs.\ trust level), and (c)~integrating the
advisor with the lifecycle manager for phase-aware
optimisation.

% ============================================================
%  Bibliography
% ============================================================
\begin{thebibliography}{25}

\bibitem{xu2008satzilla}
L.~Xu, F.~Hutter, H.~Hoos, and K.~Leyton-Brown,
``SATzilla: portfolio-based algorithm selection for SAT,''
\emph{JAIR}, vol.~32, pp.~565--606, 2008.

\bibitem{pnueli1998translation}
A.~Pnueli, M.~Siegel, and E.~Singerman,
``Translation validation,''
\emph{TACAS}, 1998.

\bibitem{chen2023codeoptimize}
M.~Chen \etal{},
``Evaluating large language models trained on code,''
\emph{arXiv preprint arXiv:2107.03374}, 2021.

\bibitem{basin1993rippling}
D.~Basin and T.~Walsh,
``A calculus for and termination of rippling,''
\emph{Journal of Automated Reasoning}, vol.~16,
pp.~147--180, 1996.

\bibitem{claessen2000quickcheck}
K.~Claessen and J.~Hughes,
``QuickCheck: a lightweight tool for random testing of
Haskell programs,''
\emph{ICFP}, 2000.

\bibitem{bansal2019holist}
K.~Bansal, S.~Loos, M.~Rabe, C.~Szegedy, and S.~Wilcox,
``HOList: an environment for machine learning of higher-order
logic theorem proving,''
\emph{ICML}, 2019.

\bibitem{hutter2011sequential}
F.~Hutter, H.~Hoos, and K.~Leyton-Brown,
``Sequential model-based optimization for general algorithm
configuration,''
\emph{LION}, 2011.

\bibitem{kotthoff2016algorithm}
L.~Kotthoff,
``Algorithm selection for combinatorial search problems:
a survey,''
\emph{AI Magazine}, vol.~35, no.~3, pp.~48--60, 2014.

\bibitem{young2024jugeo}
H.~Young,
``Judgment geometry: a framework for trust-calibrated
verification,''
\emph{JuGeo Technical Reports}, Paper~1, 2024.

\bibitem{young2024channels}
H.~Young,
``Evidence channels in judgment geometry,''
\emph{JuGeo Technical Reports}, Paper~88, 2024.

\bibitem{young2024advisors}
H.~Young,
``Copilot advisors for heap, scope, import, and callable
analysis,''
\emph{JuGeo Technical Reports}, Papers~80--81, 2024.

\bibitem{young2024research}
H.~Young,
``Copilot research and experiment advisors,''
\emph{JuGeo Technical Reports}, Paper~82, 2024.

\bibitem{young2024lifecycle}
H.~Young,
``Kernel lifecycle and health monitoring in judgment geometry,''
\emph{JuGeo Technical Reports}, Paper~93, 2024.

\bibitem{demoura2008z3}
L.~de~Moura and N.~Bj{\o}rner,
``Z3: an efficient SMT solver,''
\emph{TACAS}, 2008.

\bibitem{moura2021lean4}
L.~de~Moura and S.~Ullrich,
``The Lean~4 theorem prover and programming language,''
\emph{CADE}, 2021.

\bibitem{polu2020gptf}
S.~Polu and I.~Sutskever,
``Generative language modeling for automated theorem proving,''
\emph{arXiv preprint arXiv:2009.03393}, 2020.

\bibitem{first2022diversity}
E.~First, M.~Rabe, T.~Ringer, and Y.~Brun,
``Diversity-driven automated formal verification,''
\emph{ICSE}, 2022.

\bibitem{jiang2022thor}
A.~Q.~Jiang, W.~Li, S.~Tworber, and Y.~Wu,
``Thor: wielding hammers to integrate language models and
automated theorem provers,''
\emph{NeurIPS}, 2022.

\bibitem{lample2022hypertree}
G.~Lample \etal{},
``HyperTree proof search for neural theorem proving,''
\emph{NeurIPS}, 2022.

\bibitem{yang2019learning}
K.~Yang and J.~Deng,
``Learning to prove theorems via interacting with proof
assistants,''
\emph{ICML}, 2019.

\end{thebibliography}

% ============================================================
%  Appendix
% ============================================================
\appendix

\section{Auxiliary Definitions and Proofs}\label{app:proofs}

\begin{lemma}[Rewriting preserves obligation count bound]%
\label{lem:rewritebound}
If $o \Rightarrow O'$ and $|O'| \leq b$ for branching
factor~$b$, then after $d$ rewriting steps the total
number of obligations is at most $b^d$.
\end{lemma}
\begin{proof}
By induction on~$d$.  The base case $d = 0$ gives one
obligation.  For the inductive step, each of the at most
$b^{d-1}$ obligations rewrites into at most $b$ sub-obligations,
giving at most $b^d$ total.
\end{proof}

\begin{lemma}[Optimisation composition]%
\label{lem:optcomp}
If $\rho_1$ and $\rho_2$ are correctness-preserving
optimisations, then $\rho_2 \circ \rho_1$ is
correctness-preserving.
\end{lemma}
\begin{proof}
For any obligation~$o$:
\begin{align*}
  \mathrm{valid}(o)
  &\iff \mathrm{valid}(\rho_1(o))
    && \text{(by correctness of $\rho_1$)} \\
  &\iff \mathrm{valid}(\rho_2(\rho_1(o)))
    && \text{(by correctness of $\rho_2$)} \\
  &= \mathrm{valid}((\rho_2 \circ \rho_1)(o)).
\end{align*}
\end{proof}

\begin{lemma}[Optimisation identity]%
\label{lem:optid}
The identity function $\mathrm{id}$ is a
correctness-preserving optimisation.
\end{lemma}
\begin{proof}
$\mathrm{valid}(o) \iff \mathrm{valid}(\mathrm{id}(o))$
since $\mathrm{id}(o) = o$.
\end{proof}

\begin{corollary}[Optimisation monoid]%
\label{cor:optmonoid}
The set of correctness-preserving optimisations forms a
monoid under composition, with $\mathrm{id}$ as the identity.
\end{corollary}
\begin{proof}
Closure: Lemma~\ref{lem:optcomp}.
Identity: Lemma~\ref{lem:optid}.
Associativity: function composition is associative.
\end{proof}

\begin{lemma}[Performance model calibration]%
\label{lem:perfcalib}
If the monitor collects at least $n \geq 30$ data points,
then the performance model's mean prediction error is
bounded by $\epsilon = 2 \sigma / \sqrt{n}$ with 95\%
confidence, where $\sigma$ is the standard deviation
of observed latencies.
\end{lemma}
\begin{proof}
By the central limit theorem, the sample mean of $n$
i.i.d.\ latency observations has standard error
$\sigma / \sqrt{n}$.  The 95\% confidence interval
for the mean is $\bar{x} \pm 1.96 \sigma / \sqrt{n}$.
Since $1.96 < 2$, the mean prediction error is bounded
by $2 \sigma / \sqrt{n}$.
\end{proof}

\begin{lemma}[Non-regression checker termination]%
\label{lem:nonregterm}
The non-regression checker terminates for every pair
of obligations, provided the solver has a finite timeout.
\end{lemma}
\begin{proof}
Each verification call to the solver is bounded by the
deadline parameter.  If the solver does not respond within
the deadline, the checker records a timeout and returns
``unknown.''  Therefore the checker makes at most two
bounded calls and always terminates.
\end{proof}

\begin{lemma}[Batch validation is sound]%
\label{lem:batchsound}
If \texttt{batch\_check} returns \texttt{all\_accept = True},
then every tested pair satisfies non-regression.
\end{lemma}
\begin{proof}
\texttt{batch\_check} iterates over all pairs and returns
\texttt{all\_accept = True} only if every individual check
returns ``accept.''  By Theorem~\ref{thm:optsound}, each
accepted pair is correctness-preserving.  Therefore all
tested pairs satisfy non-regression.
\end{proof}

\begin{lemma}[Health indicator actionability]%
\label{lem:healthaction}
A \texttt{HealthIndicator} with status
\texttt{DEGRADED} or \texttt{UNHEALTHY} is actionable.
\end{lemma}
\begin{proof}
By the implementation in Listing~\ref{lst:opthealth},
the \texttt{is\_actionable()} method returns \texttt{True}
when the status is not \texttt{HEALTHY}.  \texttt{DEGRADED}
and \texttt{UNHEALTHY} are both not \texttt{HEALTHY}, hence
actionable.
\end{proof}

\begin{proposition}[Optimisation advisor safety]%
\label{prop:optsafety}
The full optimisation pipeline (suggest, validate, apply)
never installs an unsound optimisation, provided the solver
is sound and the non-regression checker is executed.
\end{proposition}
\begin{proof}
Suppose an optimisation~$\rho$ is suggested by the Copilot
channel.  Before application, the non-regression checker
verifies that $\mathrm{verify}(o) = \mathrm{verify}(\rho(o))$
using the solver.  If the solver is sound, then equal verdicts
imply equal validity.  The optimisation is applied only if the
checker accepts.  Therefore, no unsound optimisation is applied.
\end{proof}

\begin{lemma}[Rate-limit compliance for optimisation]%
\label{lem:optratelimit}
The optimisation advisor respects the global rate limit
enforced by \texttt{APIRateLimiter}.
\end{lemma}
\begin{proof}
All API calls go through \texttt{JuGeoAPI}, which enforces
the rate limit.  The advisor does not bypass the API.
Therefore, the advisor's request rate is bounded by the
global rate limit.
\end{proof}

\begin{lemma}[Trust does not accumulate through repeated optimisation]%
\label{lem:trustnoaccum}
Applying $k$ successive optimisations, each suggested at
$\Tcopilot$, does not elevate the overall trust above
$\Tcopilot$ for the suggestions.
\end{lemma}
\begin{proof}
Each optimisation suggestion has trust $\tleq \Tcopilot$
by Theorem~\ref{thm:trustpres}.  Composing $k$ suggestions
yields trust $\min(\Tcopilot, \ldots, \Tcopilot) = \Tcopilot$.
The validated evidence (from the solver) has trust $\Tsolver$,
but this is independent of the suggestion trust.
\end{proof}

\begin{remark}[Extension to multi-objective optimisation]
The strategy space can be extended with a Pareto
ordering over multiple objectives (time, trust, cost).
The ranking function then selects from the Pareto front.
We leave the formal treatment to future work.
\end{remark}

\begin{lemma}[Optimisation advisor is stateless per invocation]%
\label{lem:optstateless}
Each invocation of the optimisation advisor depends only on the
current obligation, the rewrite history, and the monitor data.
No hidden state is carried between invocations.
\end{lemma}
\begin{proof}
The advisor reads the obligation and history from its arguments,
queries the monitor for statistics, and queries the Copilot channel.
The Copilot channel maintains no conversation state between calls.
Therefore each invocation is independent.
\end{proof}

\begin{lemma}[Search-space reduction is composable]%
\label{lem:searchcomp}
If optimisation $\rho_1$ reduces the obligation count from $n$ to $m$,
and $\rho_2$ reduces from $m$ to $k$, then $\rho_2 \circ \rho_1$
reduces from $n$ to $k$.
\end{lemma}
\begin{proof}
$\rho_1$ maps $n$ obligations to $m$, and $\rho_2$ maps $m$ to $k$.
The composition $\rho_2 \circ \rho_1$ maps $n$ to $k$ directly.
The speedup factor is $n/k = (n/m) \cdot (m/k)$, the product of the
individual speedups.
\end{proof}

\begin{proposition}[Complete optimisation pipeline safety]%
\label{prop:optpipelinesafe}
The complete optimisation pipeline (suggest strategy, predict
performance, validate non-regression, apply) is safe: no
unsound optimisation is applied, trust is preserved, and
the pipeline terminates.
\end{proposition}
\begin{proof}
Safety: Theorem~\ref{thm:optsound} and Proposition~\ref{prop:optsafety}.
Trust: Theorem~\ref{thm:trustpres}.
Termination: Lemma~\ref{lem:nonregterm} (checker terminates) and
the finite number of obligations (at most one optimisation per
obligation).
\end{proof}

\begin{remark}[Interaction with lowering hints]
The optimisation advisor operates at the proof-obligation level,
before lowering to IR.  Its suggestions may affect which lowering
hints are applicable.  The two advisors compose: optimisation
rewrites are applied first, then lowering hints operate on the
optimised obligation.  Trust is propagated correctly since both
advisors are bounded by $\Tcopilot$.
\end{remark}

\end{document}
